\rok{Новембар 2024.}{E}

Први практични тест из Алгоритама и структура података

\zadatak{1}{Бинарно претраживање}
Написати sort методу класе \textit{BinarySearch}.

\textbf{Додатне напомене за решавање задатка:}
\begin{itemize}
    \item Улаз је сортиран низ целих бројева.
    \item Излаз је true или false, у зависности да ли је број пронађен.
    \item У шаблону се налази класа \texttt{SortedArrayGenerator} коју треба искористити за генерисање сортираног низа случајних бројева.
    \item Креирати класу \texttt{BinarySearch} и у оквиру ње генерисати методу:
    \begin{Verbatim}[fontsize=\footnotesize, numbers=left, xleftmargin=10mm]
public static bool search(long[] arr, int size, long value)
    \end{Verbatim}
    \item Тестирати у оквиру \texttt{Main} методе.
\end{itemize}



\zadatak{2}{Брисање елемената реда дељивих са $m$, сума преосталих}
Имплементирати методу која прима низ целих бројева и један цео број $m$. Метода треба да користи структуру података ред (класа \texttt{Queue}) за организацију и манипулацију елементима, и треба да изврши следеће кораке:

\begin{enumerate}
    \item \textbf{Уклањање бројева дељивих са m:} Проласком кроз ред, уклонити све бројеве који су дељиви са задатим бројем m.
    \item \textbf{Рачунање суме преосталих бројева:} Израчунати збир свих бројева који су остали у реду након уклањања бројева дељивих са m.
    \item \textbf{Враћање резултата:} Као резултат, метода треба да врати израчунату суму преосталих бројева у реду и исписати те бројеве.
\end{enumerate}

\textbf{Додатни захтеви за решавање задатка:}
\begin{itemize}
    \item Задатак решавати помоћу структуре реда (класа \texttt{Queue}) која је дата у оквиру шаблона задатка.
    \item Након што се уклоне бројеви дељиви са m, ред треба да остане у истом редоследу као пре уклањања тих бројева. То значи да сви преостали бројеви (они који нису дељиви са m) треба да буду враћени у ред у оригиналном редоследу.
    \item Задатак решавати у оквиру методе \texttt{public int RemoveDivisibleAndSum(int[] array, int divisor)}.
    \item Алгоритам је неопходно написати са линеарном временском комплексношћу $O(n)$.
\end{itemize}

\textbf{Примери:}
\begin{itemize}
    \item \textbf{Улаз:} \texttt{[ 3, 5, 7, 8, 10, 13, 14]}, m=3 \\
          \textbf{Излаз:} Сума елемената је 57. Преостали елементи реда: \texttt{5, 7, 8, 10, 13, 14}
    \item \textbf{Улаз:} \texttt{[12, 18, 25, 40, 9, 30, 45]}, m=5 \\
          \textbf{Излаз:} Сума елемената је 39. Преостали елементи реда: \texttt{12, 18, 9}
\end{itemize}



\zadatak{3}{Минимална и максимална удаљеност критичних тачака}
Критична тачка у повезаној листи се дефинише као локални максимум или локални минимум.

\begin{itemize}
    \item Чвор је \textbf{локални максимум} ако је вредност тренутног чвора строго већа од вредности претходног и следећег чвора.
    \item Чвор је \textbf{локални минимум} ако је вредност тренутног чвора строго мања од вредности претходног и следећег чвора.
\end{itemize}

Напомена: Чвор може бити локални максимум или минимум само ако постоји претходни и следећи чвор.

Задатак је да, с обзиром на почетни чвор повезане листе (head), вратимо низ дужине 2 који садржи:

\begin{itemize}
    \item \textbf{minDistance}: минималну удаљеност између било која два различита критична чвора
    \item \textbf{maxDistance}: максималну удаљеност између било која два различита критична чвора
\end{itemize}

Ако постоји мање од два критична чвора, треба да се врати низ [-1, -1].

\textbf{Додатни захтеви за решавање задатка:}
\begin{itemize}
    \item Задатак решавати помоћу структуре листе (класа \texttt{LinkedList}) која је дата у оквиру шаблона задатка.
    \item Не користити додатне колекције попут листе или низа, користити само повезану листу.
    \item Задатак решавати у оквиру методе \texttt{public int[] findMinMaxDistance()}.
\end{itemize}

\textbf{Примери:}
\begin{itemize}
    \item \textbf{Улаз 1:} $3 \rightarrow 1$ \\
          \textbf{Излаз 1:} \texttt{[-1,-1]}
    \item \textbf{Улаз 2:} $5 \rightarrow 3 \rightarrow 1 \rightarrow 2 \rightarrow 5 \rightarrow 1 \rightarrow 2$ \\
          \textbf{Излаз 2:} \texttt{[1,3]}
\end{itemize}

\textbf{Објашњење:}
\begin{itemize}
    \item Трећи чвор је локални минимум јер је вредност 1 мања од 3 и 2.
    \item Пети чвор је локални максимум јер је вредност 5 већа од 2 и 1.
    \item Шести чвор је локални минимум јер је вредност 1 мања од 5 и 2.
\end{itemize}

Минимална удаљеност између критичних тачака је између петог и шестог чвора, што је удаљеност \textbf{1} (6 - 5 = 1). Максимална удаљеност између критичних тачака је између трећег и шестог чвора, што је удаљеност \textbf{3} (6 - 3 = 3).

\sablon{AISP2024_Test1_GrupaE}{2024/Novembar/GrupaE/AISP2024_Test1_GrupaE.linq}
